\chapter{The Under-the-Hood ML Engine}
\label{ch:engine}

Chapter~\ref{ch:foundations} established that this project only ever runs \emph{inference}: the
weights were trained once, by someone else, and shipped inside the MediaPipe wheel. This chapter
opens the box. It explains the two-stage architecture that makes real-time hand tracking possible on
a CPU, what the model actually outputs, and the three dependency pins that will break your
installation if you ignore them.

\section{The two-stage pipeline}

A naive design would run one enormous neural network over the entire $640 \times 480$ frame, every
frame, looking for 21 landmarks. That design would be far too slow: most of the image is background,
and re-searching the whole frame 30 times a second wastes almost all of the computation.

MediaPipe's solution is to split the problem into a cheap ``where'' and an expensive ``what''.

\begin{mlconcept}{Stage 1 --- BlazePalm: find the hand, coarsely}
\textbf{BlazePalm} is a \emph{single-shot detector}. Its job is not to find fingers; it is to answer
one question, roughly: \emph{is there a palm here, and where?}

It works by sliding a set of predefined boxes --- called \textbf{anchors} --- over a downsampled
version of the frame, and for each anchor predicting two things:

\begin{enumerate}
  \item a \textbf{score}: how likely is it that this box contains a palm?
  \item an \textbf{offset}: how should the box be nudged and resized to fit better?
\end{enumerate}

It predicts a palm rather than a whole hand on purpose. A palm is rigid: its shape and size vary
little between people and poses, and --- critically --- it looks similar whether the hand is open,
closed, or rotated. A fist and an open hand have nearly the same palm box. That robustness is what
makes the detector reliable enough to hand off to stage two.

The output of stage 1 is one or more rectangles (plus a rotation angle, for hands tilted in the
image plane) that are \emph{likely} to contain a palm.
\end{mlconcept}

\begin{mlconcept}{Stage 2 --- the landmark model: find the 21 points, precisely}
Stage 2 takes the box from stage 1, \textbf{crops the frame to that box}, resizes the crop to the
network's fixed input size, and runs a second, more expensive network whose only job is regression on
21 specific points.

Because the crop is small and tightly centred on the hand, the network sees the fingers at high
effective resolution --- a fingertip may occupy 30 pixels in the crop where it occupied 5 in the full
frame. That is how it can predict a knuckle to a fraction of a pixel.

The output is 21 landmarks, each with three coordinates $(x, y, z)$ normalised to $[0, 1]$ relative
to the \emph{full} frame --- the library maps the crop coordinates back for you. The $z$ value is
relative depth, roughly ``how far in front of or behind the wrist plane'', and it is what makes this
a 3D system: fingers pointing at the camera still separate correctly in depth.
\end{mlconcept}

\begin{center}
\begin{tikzpicture}[
  box/.style={draw=slateline, fill=slatebg, rounded corners=1.2mm, align=center,
              inner sep=3mm, font=\small\sffamily, text=navy},
  lbl/.style={font=\scriptsize\sffamily, text=slatetext},
  ar/.style={-{Stealth[length=2.2mm]}, draw=cyanx, line width=1pt}
]
  % Full frame
  \draw[draw=slateline, fill=white, line width=1pt] (0,0) rectangle (3.2,2.4);
  \node[lbl] at (1.6,-0.24) {full frame $640\times480$ (downsampled)};
  \node[lbl, text=cyanx] at (1.6, 2.05) {stage 1: BlazePalm};

  % palm box inside frame
  \draw[draw=amberx, line width=1.2pt, fill=ambersoft]
        (1.0,0.55) rectangle (2.15,1.85);
  \node[lbl, text=amberx] at (1.57, 0.32) {palm box};

  % arrow to crop
  \draw[ar] (3.32,1.2) -- (4.18,1.2) node[midway, above, lbl] {crop};

  % Crop box
  \draw[draw=slateline, fill=white, line width=1pt] (4.3,0.35) rectangle (6.1,2.05);
  \node[lbl] at (5.2,-0.24) {crop, resized};
  \node[lbl, text=tealx] at (5.2, 2.05) {stage 2};

  % hand inside the crop
  \begin{scope}[shift={(5.02,0.75)}, scale=0.30]
    \handcoordinates
    \handbones
    \foreach \i in {0,...,20} {\node[joint, minimum size=4pt] at (L\i) {};}
    \foreach \i in {4,12,16,20} {\node[jointtip, minimum size=5pt] at (L\i) {};}
  \end{scope}

  % arrow to output
  \draw[ar] (6.22,1.2) -- (7.08,1.2) node[midway, above, lbl] {regress};

  % Output box
  \node[box, text width=2.5cm] (out) at (8.5,1.2)
    {21 landmarks\\ $(x,y,z)$\\ normalised};
  \node[lbl, text=tealx] at (8.5, 2.05) {stage 2 output};
\end{tikzpicture}
\captionof{figure}{The two-stage detection pipeline. Stage 1 searches the whole frame for a palm box
--- cheap, because the image is downsampled and the target is rigid. Stage 2 runs a heavier network
only inside that box. The split is what buys 30 FPS on a CPU.}
\label{fig:twostage}
\end{center}

\begin{keyidea}
\textbf{The economic argument for two stages.} Suppose the full frame costs 100 units of compute and
a crop costs 10. Running the expensive network full-frame costs 100 per frame. Running a cheap
detector once and the expensive network on a crop costs roughly $5 + 10 = 15$ --- and when the hand
is moving smoothly, MediaPipe can even skip the detector entirely for several frames and simply
\emph{track} the box, because it also predicts the hand's position in the next frame. That is the
difference between 8 FPS and 30 FPS, and it is why the design works at all.
\end{keyidea}

\section{What the model reports: the 21 landmarks}

MediaPipe numbers the 21 landmarks in a fixed order. \textbf{These indices are the vocabulary of the
entire project} --- every rule in Chapter~\ref{ch:math} is expressed as arithmetic on them, so they
are worth learning properly.

\begin{figure}[h]
\centering
\begin{tikzpicture}[scale=1.35, every node/.style={font=\scriptsize\sffamily}]
  \handcoordinates
  \handbones
  \foreach \i in {0,...,20} {\node[joint] at (L\i) {};}
  \foreach \i in {4,8,12,16,20} {\node[jointtip] at (L\i) {};}

  % emphasised landmarks
  \node[draw=redx, circle, inner sep=1.6pt, line width=0.9pt] at (L8) {};

  % labels
  \node[right=1pt, text=navy, font=\scriptsize\sffamily\bfseries] at (L0)  {0 wrist};
  \node[right=1pt, text=navy, font=\scriptsize\sffamily\bfseries] at (L4)  {4 thumb tip};
  \node[left=1pt,  text=navy, font=\scriptsize\sffamily\bfseries] at (L5)  {5 index MCP};
  \node[above right=1pt, text=navy, font=\scriptsize\sffamily\bfseries] at (L8) {8 index tip};
  \node[left=1pt,  text=navy, font=\scriptsize\sffamily\bfseries] at (L17) {17 pinky MCP};

  \node[right=1pt, text=slatetext] at (L12) {12};
  \node[left=1pt,  text=slatetext] at (L16) {16};
  \node[left=1pt,  text=slatetext] at (L20) {20};
  \node[right=1pt, text=slatetext] at (L6)  {6};
  \node[left=1pt,  text=slatetext] at (L10) {10};
  \node[left=1pt,  text=slatetext] at (L14) {14};
  \node[left=1pt,  text=slatetext] at (L18) {18};

  % annotations
  \draw[decorate, decoration={brace, amplitude=3pt, mirror}, draw=tealx]
        (1.35,0.95) -- (1.35,3.55) node[midway, right=4pt, text=tealx, align=left]
        {four long\\fingers};
\end{tikzpicture}
\caption{The MediaPipe hand topology. The bold labels are the landmarks this project reasons about
directly: the wrist anchors the palm normal, the thumb tip and IP joint drive the thumb rule, and
indices 5 and 17 define the palm plane. Tip landmarks are drawn in amber; the index tip --- the one
carrying the smoothed marker --- is ringed in red.}
\label{fig:landmarks}
\end{figure}

\subsection{The landmark table}

The naming convention is systematic, and worth internalising: \textbf{MCP} is the knuckle
(metacarpophalangeal joint), \textbf{PIP} the first finger joint, \textbf{DIP} the second, and
\textbf{TIP} the end of the finger. \textbf{CMC} is the base of the thumb, at the wrist.

\begin{center}
\rowcolors{2}{slatebg}{white}
\begin{tabularx}{\textwidth}{@{}r l l X@{}}
\toprule
\rowcolor{white}
\textbf{Index} & \textbf{Name} & \textbf{Short} & \textbf{Role in this project} \\
\midrule
0  & Wrist              & WRIST      & Origin of both palm vectors; defines the palm plane \\
1  & Thumb CMC          & ---        & \textit{(not used directly)} \\
2  & Thumb MCP          & ---        & \textit{(not used directly)} \\
3  & Thumb IP           & THUMB\_IP  & The thumb rule's ``comparison'' point \\
4  & Thumb tip          & THUMB\_TIP & The thumb rule's ``extended?'' point \\
5  & Index MCP          & INDEX\_MCP & Second palm vector endpoint; index-finger base \\
6  & Index PIP          & ---        & The knuckle index 8 must beat vertically \\
7  & Index DIP          & ---        & \textit{(not used directly)} \\
8  & Index tip          & INDEX\_TIP & The accentuated, EMA-smoothed marker \\
9  & Middle MCP         & ---        & \textit{(not used directly)} \\
10 & Middle PIP         & ---        & The knuckle index 12 must beat vertically \\
11 & Middle DIP         & ---        & \textit{(not used directly)} \\
12 & Middle tip         & ---        & Middle finger extension test \\
13 & Ring MCP           & ---        & \textit{(not used directly)} \\
14 & Ring PIP           & ---        & The knuckle index 16 must beat vertically \\
15 & Ring DIP           & ---        & \textit{(not used directly)} \\
16 & Ring tip           & ---        & Ring finger extension test \\
17 & Pinky MCP          & PINKY\_MCP & The thumb rule's fixed reference point \\
18 & Pinky PIP          & ---        & The knuckle index 20 must beat vertically \\
19 & Pinky DIP          & ---        & \textit{(not used directly)} \\
20 & Pinky tip          & ---        & Pinky finger extension test \\
\bottomrule
\end{tabularx}
\end{center}

\begin{keyidea}
\textbf{Only nine of the twenty-one landmarks do any work.} The project reads indices
$0, 3, 4, 5, 8, 12, 16, 17, 20$ --- plus the four PIP joints $6, 10, 14, 18$ as vertical references.
That is a feature, not an oversight: a rule that uses fewer inputs is easier to reason about, easier
to test, and less exposed to jitter in landmarks that carry no information.
\end{keyidea}

\subsection{Why 21 points, and not a bounding box or a mask}

\begin{mlconcept}{The representation choice, and what it buys you}
There are three common ways to represent a detected hand:

\begin{itemize}
  \item \textbf{A bounding box} --- four numbers. Cheap, but useless for fingers: a fist and an open
        hand have the same box.
  \item \textbf{A segmentation mask} --- a pixel-per-pixel ``is this hand?'' image. Precise, but
        expensive, and it still does not tell you \emph{which} pixel is a knuckle.
  \item \textbf{Keypoints} --- a small, fixed set of semantically named points. Compact, cheap to
        compare, and every point has a \emph{meaning}: index 8 \emph{is} the index fingertip.
\end{itemize}

The project needs semantics, not pixels. Its questions are ``is this fingertip above that knuckle?''
and ``how far is the thumb from the pinky knuckle?'' --- questions about \emph{identified joints}.
Keypoints answer them directly with 63 numbers ($21 \times 3$) instead of a mask with 307{,}200.
\end{mlconcept}

\section{The dependency matrix, and three package traps}
\label{sec:deps}

The project pins exactly three packages:

\begin{center}
\rowcolors{2}{slatebg}{white}
\begin{tabularx}{\textwidth}{@{}l l X@{}}
\toprule
\rowcolor{white}
\textbf{Pin} & \textbf{Installs} & \textbf{Why exactly this} \\
\midrule
\code{mediapipe==0.10.14} & the model + runtime &
  The last release family with the legacy \code{mp.solutions.hands} API the project is built on.
  Everything from 0.10.30 onward has removed it. \\
\code{opencv-contrib-python}\newline\code{==4.11.0.86} & \code{cv2} &
  The \emph{contrib} build, not \code{opencv-python}: MediaPipe already depends on it, so pinning
  the same package avoids two competing \code{cv2} installs. \\
\code{numpy<2} & numerical arrays &
  MediaPipe 0.10.x wheels are compiled against the NumPy 1.x ABI. NumPy 2.0 changed that ABI. \\
\bottomrule
\end{tabularx}
\end{center}

Each of the three is a real trap that has cost real people real hours. Take them one at a time.

\begin{edgetrap}{Trap 1 --- two \code{cv2} packages fighting over one name}
OpenCV ships under several PyPI names: \code{opencv-python}, \code{opencv-contrib-python},
\code{opencv-python-headless}, and so on. They all install the \textbf{same importable module name},
\code{cv2}, into \code{site-packages}.

\textbf{MediaPipe's own dependency list includes \code{opencv-contrib-python}.} So this
\code{requirements.txt} would be actively harmful:

\begin{lstlisting}
mediapipe==0.10.14
opencv-python==4.10.0.84      # <-- WRONG: a second, competing cv2
numpy<2
\end{lstlisting}

pip will happily install both. Then \code{import cv2} resolves to whichever package wrote its files
last, and upgrading one can silently break the other --- producing the notorious
``it works on my machine'' failures that are almost impossible to debug because
\code{cv2.\_\_version\_\_} can report one thing while the code that runs comes from the other.

\textbf{The fix.} Pin the package MediaPipe already wants. The contrib build is a strict
\emph{superset} of the base build --- it adds extra modules --- so nothing is lost:

\begin{lstlisting}
mediapipe==0.10.14
opencv-contrib-python==4.11.0.86
numpy<2
\end{lstlisting}
\end{edgetrap}

\begin{edgetrap}{Trap 2 --- NumPy 2.0 and the ABI break}
NumPy 2.0 (released mid-2024) changed its C application binary interface. Extension modules compiled
against NumPy 1.x --- which includes the MediaPipe 0.10.x wheels --- can fail at import time with
messages about binary incompatibility, or worse, import and then crash.

Pinning \code{numpy<2} resolves it to the 1.26.x line, which is also the last version supporting
Python 3.9--3.12 across the board. If you see \code{pip} proposing to upgrade NumPy as a
``helpful'' side effect of another install, decline.
\end{edgetrap}

\begin{edgetrap}{Trap 3 --- \code{mp.solutions} no longer exists (the load-bearing pin)}
This is the one that matters most, because it is not a subtlety --- it is a hard removal.

The legacy API used by this project, \code{mp.solutions.hands.Hands}, was deprecated and then
\textbf{removed} in MediaPipe 0.10.30+. On a newer version the very first attribute access fails:

\begin{lstlisting}
>>> import mediapipe as mp
>>> mp.solutions.hands
AttributeError: module 'mediapipe' has no attribute 'solutions'
\end{lstlisting}

\textbf{Consequence:} \code{pip install -U mediapipe} bricks the application. The pin
\code{mediapipe==0.10.14} is therefore not a preference --- it is a load-bearing structural element,
in the same sense as a beam holding up a floor.

\textbf{Defence.} Rather than let the user discover this through an opaque
\code{AttributeError} three imports deep, the application checks at import time and fails loudly with
an actionable message. Here is the real guard, from \code{hand\_tracker.py}:

\begin{codecard}
\begin{lstlisting}
if not hasattr(mp, "solutions") or not hasattr(mp.solutions, "hands"):
    raise RuntimeError(
        "mediapipe>=0.10.30 removed mp.solutions, which this app is built on. "
        "Install the pinned version instead: pip install -r requirements.txt "
        "(do NOT `pip install -U mediapipe`)."
    )
\end{lstlisting}
\end{codecard}

\begin{itemize}
  \item \code{hasattr(mp, "solutions")} probes the attribute \emph{without} raising --- the whole
        point of \code{hasattr} is that it returns \code{False} instead of blowing up. This is
        exactly the situation it was designed for.
  \item The \code{or} short-circuits: if \code{solutions} is missing entirely, the second
        \code{hasattr} is never evaluated, so the guard itself cannot crash.
  \item \code{raise RuntimeError} rather than \code{sys.exit} keeps the failure inside Python's
        normal error machinery, so a test harness or an editor sees a real traceback.
  \item The message names \emph{the fix}, not just the problem. A good error message tells you what
        to type next.
\end{itemize}
\end{edgetrap}

\subsection{Suppressing the deprecation noise}

Deprecated APIs like these emit warnings on import, which pollute a live console that is already
printing FPS telemetry. The project silences them \emph{only around the import}, so genuine warnings
elsewhere still surface:

\begin{codecard}
\begin{lstlisting}
with warnings.catch_warnings():
    warnings.simplefilter("ignore")
    import mediapipe as mp
\end{lstlisting}
\end{codecard}

The \code{with} block restores the previous warning behaviour on exit. Warning filters are global
state; changing them permanently to hide one import's noise would suppress diagnostics you may
actually need later --- for instance a NumPy deprecation that only appears under load.

\section{Confidence thresholds: what 0.7 and 0.6 actually control}

The tracker constructs the model with four arguments. Two of them are thresholds, and they are the
difference between a twitchy detector and a sluggish one.

\begin{codecard}
\begin{lstlisting}
self.hands = mp_hands.Hands(
    static_image_mode=mode,
    max_num_hands=max_hands,
    min_detection_confidence=det_conf,     # default 0.7
    min_tracking_confidence=track_conf,    # default 0.6
)
\end{lstlisting}
\end{codecard}

\begin{mlconcept}{Two thresholds, two different questions}
\begin{itemize}
  \item \code{min\_detection\_confidence = 0.7} applies when the system is \emph{looking for a new
        hand} --- stage 1, the palm detector running over the whole frame. A candidate box must score
        at least 0.7 to be accepted.
  \item \code{min\_tracking\_confidence = 0.6} applies when the system is \emph{following a hand it
        already found}, using the landmarks predicted from the previous frame. The bar is lower
        because tracking has more evidence: it knows where the hand just was.
\end{itemize}

Why the asymmetry? Detection is a \emph{false-positive} problem: if you accept weak candidates, a
folded sleeve or a face can spawn a phantom hand. Tracking is a \emph{false-negative} problem: if you
drop the hand for one frame, the on-screen count flickers and the debouncer restarts.

So the design is deliberately conservative about discovering hands (0.7) and forgiving about keeping
them (0.6). Lowering detection to ~0.5 is the usual first experiment for a dim room --- with the
warning that phantom hands become likelier. Raising it toward 0.9 makes the tracker refuse to see
your hand until it is well lit and centred.
\end{mlconcept}

\subsection{\code{static\_image\_mode} and \code{max\_num\_hands}}

\begin{itemize}
  \item \code{static\_image\_mode=True} disables temporal tracking. Every frame is treated as an
        unrelated still, so the detector runs every time. It is slower and noisier on video, but it
        is exactly right for \code{--image} mode, where there is only one frame and nothing to
        track. The CLI exposes it as \code{--static} for live use; \code{run\_image()} always sets
        it to \code{True}.
  \item \code{max\_num\_hands=2} caps how many hands are searched for. It is a genuine performance
        knob: the detector's cost grows with the number of candidate boxes it keeps. Two is the
        right default for a two-handed demo; raise it only if you need more.
\end{itemize}

\section{Handedness, and why the label inverts}

The model also emits a \textbf{handedness} classification --- ``Left'' or ``Right'' --- with a score.
Beginners assume this is a property of the hand in the image. It is not: it is a prediction made
under an assumption about how the image was captured.

\begin{edgetrap}{The mirrored-feed trap: MediaPipe assumes a selfie}
MediaPipe's documentation is explicit: handedness is predicted \emph{assuming the input image is
mirrored}, as when a front-facing camera feed has already been flipped horizontally --- the normal
convention for a selfie.

This project flips the frame for exactly that reason --- a mirrored view is what makes a webcam feel
like a mirror and makes ``move right'' move right on screen:

\begin{lstlisting}
frame = cv2.flip(frame, 1)
\end{lstlisting}

But that flip happens \textbf{before} MediaPipe ever sees the frame. The result: on many webcams and
configurations the label comes back \emph{inverted} --- your physical right hand is reported as
``Left''.

\textbf{Why this is handled by design rather than patched.} The project never lets the label drive
geometry. The finger logic uses handedness only in the explicitly opt-in \code{--thumb-mode x}; the
default thumb rule is handedness-free by construction (Section~\ref{sec:thumb}). The label is
treated as \emph{display telemetry}, and the escape hatch is a display-only fix:

\begin{lstlisting}
--swap-handedness     # swap the Left/Right labels in the HUD
\end{lstlisting}

Because the swap happens inside \code{parse\_hand\_data}, it also flows through to the per-hand
smoothing and debouncing keys, so no state is orphaned by the change.
\end{edgetrap}

\begin{keyidea}
\textbf{Chapter summary.} Detection is deliberately split in two: a cheap detector finds a palm box,
and an expensive network regresses 21 named 3D landmarks inside that crop --- which is what makes
real-time performance on a CPU achievable. The three pins each defend against a specific, real
failure: two competing \code{cv2} packages, the NumPy 2 ABI break, and the removal of
\code{mp.solutions}. The two confidence thresholds are tuned asymmetrically --- strict about
discovering hands, forgiving about losing them --- and the handedness label is display-only, because
on a mirrored feed it cannot be trusted.
\end{keyidea}
